% Fibonacci-Pell nearest gaps: an all-exponent classification and
% quadratic-unit orbit rigidity
%
% Authors:
%   Dr. Denys Dutykh (Mathematics Department, Khalifa University of Science
%   and Technology, Abu Dhabi, UAE)
%   Prof. Laurent Vuillon (Univ. Savoie Mont Blanc, CNRS, LAMA, Chambery,
%   France)
%
\section{Height separation in the even branch}
\label{sec:height}

Continue under the standing hypothesis of an even-exponent target with
$q\geqslant7$.
By \cref{prop:negative-exponent},

\begin{equation}
 \lambda^n-D_{q,\sigma}^2=g\lambda^{-h},
 \qquad
 h\geqslant3\ \text{odd}.
 \label{eq:negative-target}
\end{equation}

We first turn the reconstructed rectangle into a separation between $h$ and
$n$.

\begin{lemma}[Divisor-height separation]
\label{lem:height-separation}
Every target in \cref{eq:negative-target} satisfies

\begin{equation}
 \begin{array}{ll}
 \sigma=1:&h\leqslant n-5,\\
 \sigma=-1:&h\leqslant n-3.
 \end{array}
 \label{eq:height-separation}
\end{equation}

In particular, $h<n$.
\end{lemma}

\begin{proof}
Put

\begin{equation}
 L=f^2+2x^2.
\end{equation}

Since $\lambda^{-h}=-C_h+P_h\sqrt2$ for odd $h$, coefficient comparison in
\cref{eq:negative-target} gives

\begin{equation}
 gC_h=L-C_n,
 \qquad
 gP_h=P_n-2\sigma fx.
 \label{eq:coefficient-height}
\end{equation}

Let $t=g/2$. By \cref{lem:rectangle}, $t\geqslant\sqrt{k}$. For every
positive odd $h$,

\begin{equation}
 C_h=\frac{\lambda^h-\lambda^{-h}}2>\frac{\lambda^h}{3}.
 \label{eq:Ch-lower}
\end{equation}

On the plus branch, $L<A_q^2<\lambda^n$. Hence

\begin{equation}
 t=\frac{L-C_n}{2C_h}<\frac32\lambda^{n-h}.
\end{equation}

On the minus branch, \cref{eq:AB-ratio} gives

\begin{equation}
 L=\frac{A_q^2+B_q^2}{2}<5B_q^2<5\lambda^n,
\end{equation}

and therefore $t<(15/2)\lambda^{n-h}$. Combining these inequalities with
$t\geqslant\sqrt{k}$ yields

\begin{equation}
 \lambda^{2(n-h)}>
 \begin{cases}
 4k/9,&\sigma=1,\\
 4k/225,&\sigma=-1.
 \end{cases}
 \label{eq:separation-raw}
\end{equation}

The sequence $R_q$ is increasing for $q\geqslant3$: with
$f=F_q$ and $x=F_{q+1}<2f$,

\begin{equation}
 R_{q+1}-R_q=3f^2+4fx-x^2>0.
\end{equation}

Hence $k\geqslant R_7=713$, and $n-h$ is odd. Direct comparison in
\cref{eq:separation-raw} excludes $n-h\leqslant3$ in the plus case and
$n-h\leqslant1$ in the minus case. This proves
\cref{eq:height-separation}.
\end{proof}

We shall use the following explicit specialization of Matveev's theorem
\cite[Corollary~2.3]{Matveev2000}. For a positive algebraic number $\alpha$,
write $h_{\mathrm{abs}}(\alpha)$ for its absolute logarithmic Weil height.

\begin{lemma}[Matveev, three positive real logarithms]
\label{lem:matveev}
Let $K$ be a real number field of degree $D$, let
$\alpha_1,\alpha_2,\alpha_3\in K$ be positive, and let

\begin{equation}
 \Lambda=b_1\log\alpha_1+b_2\log\alpha_2+b_3\log\alpha_3\ne0
\end{equation}

with $b_i\in\Z$. If

\begin{equation}
 A_i\geqslant
 \max\{D h_{\mathrm{abs}}(\alpha_i),\abs{\log\alpha_i},0.16\}
\end{equation}

and $B\geqslant\max_i\abs{b_i}$, then

\begin{equation}
 \log\abs\Lambda>
 -1.4\mathbin{\cdot}30^6 3^{9/2}D^2(1+\log D)
 A_1A_2A_3(1+\log B).
 \label{eq:matveev-special}
\end{equation}
\end{lemma}

The first application reduces the a priori linear range for $h$ to a
logarithmic range.

\begin{proposition}[Logarithmic bound for the negative exponent]
\label{prop:h-log-bound}
Every target with $q\geqslant7$ satisfies

\begin{equation}
 h<2.4\mathbin{\cdot}10^{14}\bigl(1+\log(2q)\bigr)
 \label{eq:h-log-bound}
\end{equation}

and

\begin{equation}
 n<2q.
 \label{eq:n-less-2q}
\end{equation}
\end{proposition}

\begin{proof}
Put

\begin{equation}
 \varphi=\frac{1+\sqrt5}{2},
 \qquad
 \epsilon=(-1)^q,
 \qquad
 u=\sigma+\sqrt2\varphi,
 \qquad
 v=-\sigma+\sqrt2\varphi^{-1}.
 \label{eq:u-v}
\end{equation}

Binet's formula gives

\begin{equation}
 D_{q,\sigma}
 =\frac{u\varphi^q+\epsilon v\varphi^{-q}}{\sqrt5}.
 \label{eq:D-binet}
\end{equation}

Set

\begin{equation}
 \delta=\epsilon\frac vu\varphi^{-2q},
 \qquad
 A_0=\frac{u^2}{5}\varphi^{2q}.
\end{equation}

Then $D_{q,\sigma}^2=A_0(1+\delta)^2$. The elementary bounds
$1<u<4$ and $\abs{v/u}<2$ give

\begin{equation}
 \abs\delta<2\left(\frac23\right)^{14}<\frac1{100}.
 \label{eq:delta-first}
\end{equation}

By \cref{lem:height-separation} and the nearest window,

\begin{equation}
 \lambda^{h+3}<2\frac{16}{5}\left(\frac{101}{100}\right)^2
 \varphi^{2q}<7\varphi^{2q}.
\end{equation}

Since $\lambda>12/5$, this strengthens \cref{eq:delta-first} to

\begin{equation}
 \abs\delta<\frac{1750}{1728}\lambda^{-h}.
 \label{eq:delta-second}
\end{equation}

By \cref{lem:rectangle}, $g<k=A_qB_q$. Therefore
$g/D_{q,\sigma}^2<3$ for either sign, by \cref{eq:AB-ratio}. Define

\begin{equation}
 \Lambda_0=n\log\lambda-2q\log\varphi
 -\log\left(\frac{u^2}{5}\right).
 \label{eq:Lambda0}
\end{equation}

After dividing \cref{eq:negative-target} by $A_0$, the estimates
\cref{eq:delta-first,eq:delta-second} give

\begin{equation}
 \abs{e^{\Lambda_0}-1}<\frac{128}{25}\lambda^{-h}<\frac{10}{27}.
\end{equation}

The elementary logarithm estimate on this interval yields

\begin{equation}
 0<\abs{\Lambda_0}<9\lambda^{-h}.
 \label{eq:Lambda0-small}
\end{equation}

The nonvanishing is exact. In
$K=\Q(\sqrt2,\sqrt5)$ one has
$N_{K/\Q}(u)=-1$ and
$N_{K/\Q}(\lambda)=N_{K/\Q}(\varphi)=1$. Equality in
\cref{eq:Lambda0} would give $1=5^{-4}$ after taking full norms.

We next verify \cref{eq:n-less-2q}. From \cref{eq:D-binet},
$D_{q,\sigma}<4\varphi^q$. Also $\lambda/\varphi>7/5$ and
$(7/5)^{14}>32$. Thus, for $q\geqslant7$,

\begin{equation}
 \lambda^{2q}>32\varphi^{2q}>2D_{q,\sigma}^2>\lambda^n.
\end{equation}

For \cref{lem:matveev}, the field degree is $4$ and admissible parameters
for \cref{eq:Lambda0} are

\begin{equation}
 A_1=2\log\lambda,
 \qquad
 A_2=2\log\varphi,
 \qquad
 A_3=20,
 \qquad
 B=2q.
\end{equation}

Indeed,

\begin{equation}
 h_{\mathrm{abs}}(u^2/5)
 \leqslant3\log2+\log\varphi+\log5<5,
\end{equation}

and $1/5<u^2/5<16/5$. The constant before $A_3$ satisfies the exact
outward inequality

\begin{equation}
 1.4\mathbin{\cdot}30^6 3^{9/2}4^2(1+\log4)
 (2\log\lambda)(2\log\varphi)
 <9.275\mathbin{\cdot}10^{12}.
 \label{eq:matveev-prefactor}
\end{equation}

Writing $H=1+\log(2q)$, comparison with
\cref{eq:Lambda0-small} gives

\begin{equation}
 h\log\lambda<1.855\mathbin{\cdot}10^{14}H+\log9.
\end{equation}

Finally $\log\lambda>4/5$. Rounding only upward proves
\cref{eq:h-log-bound}.
\end{proof}

The fixed-coefficient form has therefore done two jobs: it makes $h$
logarithmic in $q$ and keeps $n$ below $2q$. It does not yet bound $q$.
The next section supplies that missing absolute bound by allowing the third
logarithmic coefficient to depend on the now-controlled exponent $h$.
